\chapter{Étude préliminaire}

\section*{Introduction}
Ce chapitre pose les fondations analytiques du projet. Nous identifions les acteurs
du système, formalisons les besoins fonctionnels et non fonctionnels, construisons
le product backlog avec \textbf{34 user stories} réparties sur \textbf{6 sprints},
puis décrivons l'architecture globale et la pile technologique retenue.

% ============================================================
\section{Les acteurs du système}

\subsection{Rôles applicatifs}

\begin{table}[H]
\centering
\renewcommand{\arraystretch}{1.5}
\begin{tabular}{|p{3cm}|p{10cm}|}
\hline
\textbf{Rôle} & \textbf{Description} \\ \hline

\textbf{Admin} &
Accès complet : gestion des utilisateurs, équipes, audit, toutes les fonctionnalités. 
Supervise l'ensemble des projets. \\ \hline

\textbf{Member} &
Profil standard : gestion de ses projets, tâches, IA, Gantt, calendrier, notifications. \\ \hline

\textbf{Client} &
Accès en lecture seule aux projets auxquels il est invité. 
Peut consulter l'avancement et les livrables. \\ \hline

\end{tabular}
\caption{Description des rôles utilisateurs}
\label{tab:roles}
\end{table}

\subsection{Rôles dans un projet}

\begin{table}[H]
  \centering
  \caption{Rôles dans un projet et droits associés}
  \label{tab:roles_projet}
  \begin{tabularx}{\textwidth}{L{2.5cm}X}
    \toprule
    \thead{Rôle projet} & \thead{Droits} \\
    \midrule
    \rowcolor{lightgray}
    Product Owner & Toutes les actions sur le projet, y compris suppression et archivage \\
    Manager & Modifier le projet, gérer les membres, créer/supprimer des tâches et statuts \\
    \rowcolor{lightgray}
    Développeur & Créer et modifier des tâches, ajouter des commentaires, saisir du temps \\
    Client & Lecture seule : consulter les tâches, le Gantt, les commentaires \\
    \bottomrule
  \end{tabularx}
\end{table}



\subsection{Acteurs logiciels}

\begin{table}[H]
\centering
\renewcommand{\arraystretch}{1.5}
\begin{tabular}{|p{4cm}|p{9cm}|}
\hline
\textbf{Composant} & \textbf{Description} \\ \hline

\textbf{Service ML} &
Microservice FastAPI (Python 3.11) exposant les endpoints de recommandation, 
prédiction de retard et estimation de durée. \\ \hline

\textbf{LLM OpenRouter} &
API tierce (Llama 3.3 70B Instruct) invoquée pour la génération 
de justifications en langage naturel. \\ \hline

\end{tabular}
\caption{Description des composants IA}
\label{tab:composantsIA}
\end{table}

\subsection{Matrice des droits}

\begin{table}[H]
  \centering
  \caption{Matrice des droits par rôle système}
  \label{tab:matrice_droits}
  \small
  \begin{tabularx}{\textwidth}{L{5cm}C{2.5cm}C{2.5cm}C{2.5cm}}
    \toprule
    \thead{Action} & \thead{Admin} & \thead{Member} & \thead{Client} \\
    \midrule
    \rowcolor{lightgray}
    Gérer les utilisateurs          & \cmark & \xmark & \xmark \\
    Gérer les équipes               & \cmark & \xmark & \xmark \\
    \rowcolor{lightgray}
    Créer un projet                 & \cmark & \cmark & \xmark \\
    Accéder au Gantt                & \cmark & \cmark & \xmark \\
    \rowcolor{lightgray}
    Accéder au calendrier           & \cmark & \cmark & \xmark \\
    Accéder aux fonctionnalités IA  & \cmark & \cmark & \xmark \\
    \rowcolor{lightgray}
    Consulter ses projets           & \cmark & \cmark & \cmark \\
    Consulter le journal d'audit    & \cmark & \cmark & \xmark \\
    \rowcolor{lightgray}
    Gérer les rôles et permissions  & \cmark & \xmark & \xmark \\
    Exporter rapport PDF d'audit    & \cmark & \xmark & \xmark \\
    \rowcolor{lightgray}
    Recevoir des notifications      & \cmark & \cmark & \cmark \\
    Modifier son profil             & \cmark & \cmark & \cmark \\
    \bottomrule
  \end{tabularx}
\end{table}

% ============================================================
\section{Besoins fonctionnels}

\begin{itemize}

    \item \textbf{Authentification et gestion des utilisateurs}
    \begin{itemize}
        \item \textbf{BF-AUTH-01 :} S'inscrire et se connecter
        \item \textbf{BF-AUTH-02 :} Réinitialiser le mot de passe par email
        \item \textbf{BF-AUTH-03 :} Gérer son profil et ses préférences
    \end{itemize}

    \item \textbf{Gestion des projets}
    \begin{itemize}
        \item \textbf{BF-PROJ-01 :} Créer et configurer un projet
        \item \textbf{BF-PROJ-02 :} Inviter des membres par email
        \item \textbf{BF-PROJ-03 :} Gérer les statuts Kanban personnalisés
    \end{itemize}

    \item \textbf{Gestion des sprints}
    \begin{itemize}
        \item \textbf{BF-SPRINT-01 :} Créer et planifier un sprint
        \item \textbf{BF-SPRINT-02 :} Démarrer et clôturer un sprint avec report automatique du backlog
    \end{itemize}

    \item \textbf{Gestion des tâches}
    \begin{itemize}
        \item \textbf{BF-TASK-01 :} Créer des tâches hiérarchiques
        \item \textbf{BF-TASK-02 :} Définir des dépendances FS/SS/FF/SF
        \item \textbf{BF-TASK-03 :} Affecter des membres avec allocation en pourcentage
        \item \textbf{BF-TASK-04 :} Déplacer les tâches en Kanban par drag-and-drop
    \end{itemize}

    \item \textbf{Planification}
    \begin{itemize}
        \item \textbf{BF-PLAN-01 :} Visualiser le Gantt avec chemin critique
        \item \textbf{BF-PLAN-02 :} Consulter le calendrier mensuel
    \end{itemize}

    \item \textbf{Notifications}
    \begin{itemize}
        \item \textbf{BF-NOTIF-01 :} Recevoir des notifications temps réel in-app
        \item \textbf{BF-NOTIF-02 :} Configurer les préférences de notification
    \end{itemize}

    \item \textbf{Dashboard et audit}
    \begin{itemize}
        \item \textbf{BF-DASH-01 :} Consulter un dashboard personnalisé selon le rôle
        \item \textbf{BF-DASH-02 :} Accéder au journal d'audit
    \end{itemize}

    \item \textbf{Fonctionnalités IA}
    \begin{itemize}
        \item \textbf{BF-IA-01 :} Obtenir une recommandation ML d'assignation
        \item \textbf{BF-IA-02 :} Estimer la durée d'une tâche par LLM
        \item \textbf{BF-IA-03 :} Détecter les risques
        \item \textbf{BF-IA-04 :} Utiliser le chat assistant IA
    \end{itemize}

    \item \textbf{Sécurité enterprise et RBAC}
    \begin{itemize}
        \item \textbf{BF-SEC-01 :} Créer et gérer des rôles personnalisés
        \item \textbf{BF-SEC-02 :} Affecter des permissions atomiques aux rôles
        \item \textbf{BF-SEC-03 :} Consulter la matrice globale des permissions
        \item \textbf{BF-SEC-04 :} Consulter le journal d'audit enrichi (valeurs avant/après, IP, sévérité)
        \item \textbf{BF-SEC-05 :} Analyser l'activité projet via le dashboard analytics (KPIs, anomalies)
        \item \textbf{BF-SEC-06 :} Exporter le journal d'audit en rapport PDF
    \end{itemize}

\end{itemize}% ============================================================
\section{Besoins non fonctionnels}

\begin{table}[H]
  \centering
  \caption{Besoins non fonctionnels}
  \label{tab:bnf}
  \begin{tabularx}{\textwidth}{L{3cm}L{5cm}X}
    \toprule
    \thead{Catégorie} & \thead{Exigence} & \thead{Valeur cible} \\
    \midrule
    \rowcolor{lightgray}
    Sécurité     & Auth stateless Bearer Token         & Sanctum ; expiration configurable \\
    Sécurité     & RBAC à 3 niveaux                   & Spatie (système) + ProjectPolicy (projet) + permissions atomiques \\
    \rowcolor{lightgray}
    Sécurité     & Journal d'audit enterprise         & Traçabilité IP, sévérité auto, valeurs avant/après \\
    \rowcolor{lightgray}
    Performance  & Latence P95 API                    & $\leq$ 500 ms \\
    Performance  & Inférence ML (P50)                 & $\leq$ 45 ms \\
    \rowcolor{lightgray}
    Ergonomie    & Responsive (breakpoints)           & Mobile, tablette, desktop \\
    Ergonomie    & Accessibilité                      & WCAG 2.1 niveau AA \\
    \rowcolor{lightgray}
    Fiabilité    & Disponibilité cible                & 99,5\,\% \\
    Fiabilité    & Détection cycle dépendances        & BFS avant toute insertion \\
    \rowcolor{lightgray}
    Maintenabilité & Duplication code (SonarCloud)    & $\leq$ 3\,\% \\
    Maintenabilité & Couverture tests (backend)       & $\geq$ 80\,\% \\
    \bottomrule
  \end{tabularx}
\end{table}

% ============================================================
\section{Product Backlog — User Stories}

\begin{longtable}{C{1.2cm}L{2cm}L{6.5cm}C{2cm}C{1.5cm}C{1.5cm}}
  \caption{Product Backlog — 34 User Stories}
  \label{tab:product_backlog}\\
  \toprule
  \thead{ID} & \thead{Acteur} & \thead{User Story} & \thead{MoSCoW} & \thead{Pts} & \thead{Sprint} \\
  \midrule
  \endfirsthead
  \multicolumn{6}{c}{\small\textit{(suite du tableau \ref{tab:product_backlog})}} \\
  \toprule
  \thead{ID} & \thead{Acteur} & \thead{User Story} & \thead{MoSCoW} & \thead{Pts} & \thead{Sprint} \\
  \midrule
  \endhead
  \bottomrule
  \endfoot

  \rowcolor{lightgray}
  US-01 & Tous    & S'inscrire et se connecter en sécurité              & Must    & 5  & S1 \\
  US-02 & Admin   & Gérer les comptes utilisateurs                      & Must    & 8  & S1 \\
  \rowcolor{lightgray}
  US-03 & Admin   & Définir rôles et permissions Spatie                  & Must    & 5  & S1 \\
  US-04 & Tous    & Modifier profil et préférences                      & Should  & 5  & S1 \\
  \rowcolor{lightgray}
  US-05 & Tous    & Réinitialiser mot de passe par email                & Should  & 3  & S1 \\
  US-06 & Member  & Créer et configurer un projet                       & Must    & 5  & S2 \\
  \rowcolor{lightgray}
  US-07 & Chef    & Inviter des membres par email (7 jours)              & Must    & 8  & S2 \\
  US-08 & Chef    & Définir des statuts Kanban personnalisés             & Should  & 5  & S2 \\
  \rowcolor{lightgray}
  US-09 & Chef    & Créer et démarrer un sprint                         & Must    & 8  & S2 \\
  US-10 & Chef    & Clôturer un sprint avec report du backlog           & Should  & 5  & S2 \\
  \rowcolor{lightgray}
  US-11 & Member  & Créer des tâches hiérarchiques (sous-tâches)        & Must    & 8  & S3 \\
  US-12 & Member  & Définir des dépendances FS/SS/FF/SF                 & Must    & 8  & S3 \\
  \rowcolor{lightgray}
  US-13 & Member  & Affecter des membres avec allocation \%             & Must    & 5  & S3 \\
  US-14 & Member  & Déplacer une tâche en Kanban (drag-and-drop)        & Must    & 5  & S3 \\
  \rowcolor{lightgray}
  US-15 & Member  & Saisir le temps passé sur une tâche                 & Should  & 3  & S3 \\
  US-16 & Member  & Commenter une tâche et mentionner (@)               & Should  & 5  & S3 \\
  \rowcolor{lightgray}
  US-17 & Member  & Visualiser le Gantt avec chemin critique            & Must    & 13 & S4 \\
  US-18 & Member  & Consulter le calendrier mensuel                     & Should  & 5  & S4 \\
  \rowcolor{lightgray}
  US-19 & Tous    & Recevoir des notifications temps réel               & Must    & 8  & S4 \\
  US-20 & Tous    & Configurer ses préférences de notification          & Should  & 3  & S4 \\
  \rowcolor{lightgray}
  US-21 & Member  & Accéder à un dashboard personnel                   & Should  & 8  & S4 \\
  US-22 & Admin   & Accéder au dashboard administrateur                 & Should  & 5  & S4 \\
  \rowcolor{lightgray}
  US-23 & Member  & Obtenir une recommandation ML d'assignation         & Must    & 13 & S5 \\
  US-24 & Member  & Estimer la durée d'une tâche par IA                 & Must    & 8  & S5 \\
  \rowcolor{lightgray}
  US-25 & Member  & Discuter avec l'assistant IA                        & Should  & 8  & S5 \\
  US-26 & Admin   & Relancer l'entraînement du modèle ML                & Could   & 5  & S5 \\
  \rowcolor{lightgray}
  US-27 & Client  & Consulter le rapport d'avancement de son projet     & Should  & 5  & S4 \\
  US-28 & Tous    & Utiliser la recherche globale                       & Should  & 5  & S4 \\
  \rowcolor{lightgray}
  US-29 & Admin   & Créer et gérer des rôles personnalisés              & Must    & 8  & S6 \\
  US-30 & Admin   & Affecter des permissions atomiques aux rôles         & Must    & 8  & S6 \\
  \rowcolor{lightgray}
  US-31 & Admin   & Consulter la matrice globale des permissions         & Should  & 5  & S6 \\
  US-32 & Admin/Member & Consulter le journal d'audit enrichi (avant/après) & Must & 8 & S6 \\
  \rowcolor{lightgray}
  US-33 & Admin   & Analyser l'activité via le dashboard analytics      & Should  & 8  & S6 \\
  US-34 & Admin   & Exporter le journal d'audit en PDF                  & Could   & 5  & S6 \\
  \midrule
  \multicolumn{4}{r}{\textbf{Total}} & \textbf{228} & \textbf{6 sprints} \\
\end{longtable}

% ============================================================
\section{Diagramme global des cas d'utilisation}

\begin{figure}[H]
  \centering
  \IfFileExists{figures/ch2_global_use_case.png}{%
    \includegraphics[width=0.92\textwidth]{figures/ch2_global_use_case.png}%
  }{%
    \fbox{\parbox{0.92\textwidth}{\centering\vspace{2cm}
      \textbf{[IMAGE A INSÉRER]}\\[0.4em]
      \texttt{figures/ch2\_global\_use\_case.png}\\[0.3em]
      \textit{Diagramme global des cas d'utilisation — 3 acteurs humains, 2 acteurs logiciels,
      5 packages fonctionnels}
    \vspace{2cm}}}%
  }
  \caption{Diagramme global des cas d'utilisation — TAC-TIC IntelliPM}
  \label{fig:global_use_case}
\end{figure}

Le diagramme regroupe les 34 cas d'utilisation en six packages :
\textbf{Authentification}, \textbf{Projets \& Équipes}, \textbf{Tâches},
\textbf{Planification}, \textbf{Intelligence Artificielle} et
\textbf{Sécurité enterprise \& Audit}. Les acteurs logiciels
(Service ML FastAPI, LLM OpenRouter, Système WebSocket) interviennent en tant
qu'acteurs secondaires dans leurs packages respectifs.

% ============================================================
\section{Conception architecturale}

\subsection{Architecture en couches}

\begin{figure}[H]
  \centering
  \IfFileExists{figures/ch2_architecture_couches.png}{%
    \includegraphics[width=0.85\textwidth]{figures/ch2_architecture_couches.png}%
  }{%
    \fbox{\parbox{0.85\textwidth}{\centering\vspace{1.5cm}
      \textbf{[IMAGE A INSÉRER]}\\[0.4em]
      \texttt{figures/ch2\_architecture\_couches.png}\\[0.3em]
      \textit{Architecture en 5 couches : Présentation, Services applicatifs, Métier,
      Persistance, IA}
    \vspace{1.5cm}}}%
  }
  \caption{Architecture en couches de TAC-TIC IntelliPM}
  \label{fig:archi_couches}
\end{figure}

\subsection{Architecture distribuée}

\begin{figure}[H]
  \centering
  \IfFileExists{figures/ch2_architecture_distribuee.png}{%
    \includegraphics[width=0.90\textwidth]{figures/ch2_architecture_distribuee.png}%
  }{%
    \fbox{\parbox{0.90\textwidth}{\centering\vspace{1.5cm}
      \textbf{[IMAGE A INSÉRER]}\\[0.4em]
      \texttt{figures/ch2\_architecture\_distribuee.png}\\[0.3em]
      \textit{Architecture distribuée : React SPA $\leftrightarrow$ Laravel API $\leftrightarrow$ MySQL + FastAPI ML}
    \vspace{1.5cm}}}%
  }
  \caption{Architecture distribuée — Vue déploiement}
  \label{fig:archi_distribuee}
\end{figure}

\subsection{Pile technologique}

\begin{table}[H]
  \centering
  \caption{Pile technologique complète}
  \label{tab:techstack}
  \small
  \begin{tabularx}{\textwidth}{L{3cm}L{3.5cm}C{1.8cm}X}
    \toprule
    \thead{Catégorie} & \thead{Outil} & \thead{Version} & \thead{Rôle} \\
    \midrule
    \rowcolor{lightgray}
    Frontend      & React                  & 19       & SPA, composants UI \\
    Frontend      & Vite                   & 7        & Bundler, HMR \\
    \rowcolor{lightgray}
    Frontend      & TailwindCSS            & 4        & Styles utilitaires \\
    Frontend      & TanStack React Query   & v5       & Cache et état serveur \\
    \rowcolor{lightgray}
    Frontend      & React Router DOM       & 6        & Routing SPA \\
    Frontend      & Framer Motion          & --       & Animations \\
    \rowcolor{lightgray}
    Backend       & Laravel                & 12       & Framework API REST \\
    Backend       & PHP                    & 8.2      & Langage serveur \\
    \rowcolor{lightgray}
    Backend       & Laravel Sanctum        & --       & Auth Bearer Token \\
    Backend       & Spatie Permission      & --       & RBAC système \\
    \rowcolor{lightgray}
    Backend       & Laravel Reverb         & --       & WebSocket server \\
    Backend       & Eloquent ORM           & --       & Mapping objet-relationnel \\
    \rowcolor{lightgray}
    Base de données & MySQL                & 8        & Persistance principale \\
    IA/ML         & FastAPI                & --       & Microservice ML \\
    \rowcolor{lightgray}
    IA/ML         & Scikit-learn           & 1.4      & GBM, RF, LogReg \\
    IA/ML         & OpenRouter (LLM)       & --       & Llama 3.3 70B Instruct \\
    \rowcolor{lightgray}
    Qualité       & SonarCloud             & --       & Analyse statique \\
    \bottomrule
  \end{tabularx}
\end{table}

\subsection{Diagramme de classes global}

\begin{figure}[H]
  \centering
  \IfFileExists{figures/ch2_class_diagram.png}{%
    \includegraphics[width=0.95\textwidth]{figures/ch2_class_diagram.png}%
  }{%
    \fbox{\parbox{0.95\textwidth}{\centering\vspace{2cm}
      \textbf{[IMAGE A INSÉRER]}\\[0.4em]
      \texttt{figures/ch2\_class\_diagram.png}\\[0.3em]
      \textit{Diagramme de classes global — 20 entités avec relations et cardinalités}
    \vspace{2cm}}}%
  }
  \caption{Diagramme de classes global de TAC-TIC IntelliPM — 20 entités}
  \label{fig:class_global}
\end{figure}

\begin{table}[H]
  \centering
  \caption{Les 20 entités du modèle de données}
  \label{tab:entites}
  \begin{tabularx}{\textwidth}{L{4cm}X}
    \toprule
    \thead{Classe} & \thead{Rôle} \\
    \midrule
    \rowcolor{lightgray}
    User            & Utilisateur du système (profil, auth, préférences) \\
    Team            & Équipe de développement (leader, membres) \\
    \rowcolor{lightgray}
    Specialty       & Compétence/technologie associée à un utilisateur \\
    Project         & Projet avec code, statut, dates, équipe \\
    \rowcolor{lightgray}
    ProjectInvitation & Invitation email (token UUID, expiration 7 jours) \\
    Sprint          & Sprint Scrum avec objectif, dates, statut \\
    \rowcolor{lightgray}
    TaskStatus      & Colonne Kanban (nom, couleur, ordre) \\
    Task            & Tâche avec hiérarchie, priorité, workflow, IA \\
    \rowcolor{lightgray}
    TaskAssignment  & Assignation avec pourcentage d'allocation \\
    TaskDependency  & Dépendance FS/SS/FF/SF entre tâches \\
    \rowcolor{lightgray}
    TaskComment     & Commentaire (threaded) avec mentions @ \\
    Attachment      & Pièce jointe (fichier ou lien) \\
    \rowcolor{lightgray}
    TimeEntry       & Temps passé (manuel ou timer) \\
    Notification    & Notification in-app avec type et données JSON \\
    \rowcolor{lightgray}
    ActivityLog     & Journal d'audit enterprise (13 champs : old/new values, IP, sévérité, actor\_role…) \\
    Role            & Rôle RBAC personnalisable (Spatie) avec permissions atomiques \\
    \rowcolor{lightgray}
    Permission      & Permission atomique (35 permissions granulaires) \\
    AiRecommendation & Résultat d'une requête IA (ML ou LLM) \\
    \rowcolor{lightgray}
    PersonalAccessToken & Token Sanctum pour l'authentification \\
    ProjectMember   & Pivot projet-utilisateur avec rôle projet (Owner/Manager/Dev/Client) \\
    \bottomrule
  \end{tabularx}
\end{table}


% ============================================================
\section{Environnement de travail et choix techniques}

\subsection{Environnement matériel}

Le développement de TAC-TIC IntelliPM a été réalisé sur les postes de travail
suivants :

\begin{table}[H]
  \centering
  \caption{Environnement matériel utilisé}
  \label{tab:env_materiel}
  \begin{tabularx}{\textwidth}{L{3.5cm}X}
    \toprule
    \thead{Composant} & \thead{Spécification} \\
    \midrule
    \rowcolor{lightgray}
    Processeur     & Intel Core i7-12\textsuperscript{e} génération / AMD Ryzen 7 5800H \\
    Mémoire RAM    & 16\,Go DDR4 \\
    \rowcolor{lightgray}
    Stockage       & SSD NVMe 512\,Go \\
    Système d'exploitation & Windows 11 Pro / Ubuntu 22.04 LTS \\
    \rowcolor{lightgray}
    Navigateur     & Google Chrome 124+ / Firefox 125+ \\
    Résolution     & 1920 × 1080 (Full HD) minimum \\
    \bottomrule
  \end{tabularx}
\end{table}

\subsection{Environnement logiciel}

\begin{table}[H]
  \centering
  \caption{Environnement logiciel de développement}
  \label{tab:env_logiciel}
  \begin{tabularx}{\textwidth}{L{4.5cm}L{3cm}X}
    \toprule
    \thead{Outil} & \thead{Version} & \thead{Usage} \\
    \midrule
    \rowcolor{lightgray}
    Visual Studio Code    & 1.89+     & Éditeur principal (backend, frontend, ML) \\
    PHP Storm             & 2024.1    & Débogage avancé Laravel \\
    \rowcolor{lightgray}
    Docker Desktop        & 4.30+     & Conteneurisation locale \\
    Postman               & 11+       & Tests API REST manuels \\
    \rowcolor{lightgray}
    TablePlus             & 5.x       & Administration base de données MySQL \\
    Git                   & 2.44+     & Contrôle de version local \\
    \rowcolor{lightgray}
    Node.js               & 22 LTS    & Runtime JavaScript / outillage Vite \\
    Composer              & 2.7+      & Gestionnaire de dépendances PHP \\
    \bottomrule
  \end{tabularx}
\end{table}

\subsection{Frameworks et technologies}

\subsubsection{Backend : Laravel 12}

\textbf{Laravel~12} (PHP~8.2) a été retenu pour le backend API REST pour plusieurs
raisons : son écosystème mature (Sanctum, Reverb, Horizon, Spatie), son ORM Eloquent
hautement expressif, et sa gestion native des queues et événements. L'architecture
suit le pattern \textbf{Repositories + Services} pour séparer la logique métier des
contrôleurs.

\subsubsection{Frontend : React 19}

\textbf{React~19} avec \textbf{Vite~7} constitue le cœur du frontend SPA. L'état
serveur est géré par \textbf{TanStack React Query v5} (cache, invalidation, pagination),
le routing par \textbf{React Router DOM v6}, les animations par \textbf{Framer Motion},
et le drag-and-drop Kanban par \textbf{@dnd-kit}. L'ensemble de l'interface est
stylisé avec \textbf{TailwindCSS~v4}.

\subsubsection{Microservice IA : FastAPI / Python}

Le microservice d'intelligence artificielle est développé en \textbf{Python~3.11}
avec le framework \textbf{FastAPI}. Les modèles de Machine Learning utilisent
\textbf{Scikit-learn~1.4} (Random Forest, Gradient Boosting, régression logistique).
Le serveur ASGI est \textbf{Uvicorn} (2 workers en production). Les modèles entraînés
sont sérialisés au format \textbf{Joblib} pour un chargement rapide en mémoire.

\subsubsection{Base de données : MySQL 8}

\textbf{MySQL~8} a été préféré aux solutions NoSQL pour la cohérence transactionnelle
requise par les relations complexes du modèle (18 entités avec clés étrangères,
contraintes d'intégrité pour les dépendances de tâches). Les migrations sont gérées
par Eloquent et les requêtes optimisées par des index sur les colonnes critiques.

\subsubsection{Qualité et analyse du code : SonarCloud}

\textbf{SonarCloud} (édition cloud de SonarQube) assure l'analyse statique continue
du code source. Il est intégré au pipeline GitHub Actions et évalue quatre dimensions :
\textbf{Reliability} (bugs), \textbf{Security} (vulnérabilités), \textbf{Maintainability}
(code smells) et \textbf{Duplication}. Le Quality Gate personnalisé exige : 0 bug critique,
0 vulnérabilité de sévérité haute, duplication $\leq$ 3\,\%, couverture $\geq$ 80\,\%.


\subsubsection{Versionnage et tests}

Le code source est hébergé sur \textbf{GitHub} (monorepo). Le workflow Git adopté est
\textbf{GitHub Flow} : branches de fonctionnalités courtes, pull requests avec review
obligatoire, merge sur \texttt{main} déclenche le pipeline CI/CD. Les tests sont
exécutés automatiquement sur chaque pull request avant toute fusion.

\subsection{Langages de programmation}

\begin{table}[H]
  \centering
  \caption{Langages de programmation utilisés}
  \label{tab:langages}
  \begin{tabularx}{\textwidth}{L{3.5cm}L{2.5cm}X}
    \toprule
    \thead{Langage} & \thead{Version} & \thead{Domaine d'utilisation} \\
    \midrule
    \rowcolor{lightgray}
    PHP           & 8.2       & Backend Laravel (contrôleurs, services, modèles) \\
    JavaScript    & ES2024    & Frontend React, hooks, composants, utilitaires \\
    \rowcolor{lightgray}
    TypeScript    & 5.x       & Typage strict des interfaces React et API \\
    Python        & 3.11      & Microservice ML (FastAPI, Scikit-learn, Pandas) \\
    \rowcolor{lightgray}
    SQL           & MySQL 8   & Requêtes Eloquent, migrations, seeders \\
    YAML          & --        & Configuration Docker Compose, GitHub Actions \\
    \rowcolor{lightgray}
    Bash          & --        & Scripts de déploiement, entrypoints Docker \\
    \bottomrule
  \end{tabularx}
\end{table}

\subsection{Base de données / SGBD}

\textbf{MySQL~8} avec le moteur \textbf{InnoDB} est le SGBD relationnel utilisé.
Les choix de configuration retenus : jeu de caractères \texttt{utf8mb4} (support
Unicode complet), niveau d'isolation \texttt{REPEATABLE READ}, et index composites
sur les colonnes les plus filtrées (\texttt{project\_id + status\_id},
\texttt{assignee\_id + due\_date}). Un cache de requêtes Redis est interposé pour
les endpoints à forte charge (dashboard, Gantt).

\subsection{Outils de conception UML}

\begin{table}[H]
  \centering
  \caption{Outils de modélisation et conception}
  \label{tab:outils_uml}
  \begin{tabularx}{\textwidth}{L{3.5cm}X}
    \toprule
    \thead{Outil} & \thead{Usage dans le projet} \\
    \midrule
    \rowcolor{lightgray}
    \textbf{Miro}         & Diagrammes UML collaboratifs : cas d'utilisation, séquence,
                            classes — tableau blanc en ligne partagé avec l'encadrant \\
    \textbf{PlantUML}     & Génération automatique de diagrammes UML depuis le code
                            (intégration VS Code) \\
    \rowcolor{lightgray}
    \textbf{Draw.io}      & Architecture distribuée, flux de déploiement, schémas réseau \\
    \textbf{Figma}        & Maquettes haute fidélité des interfaces (wireframes → UI finale) \\
    \bottomrule
  \end{tabularx}
\end{table}


\subsection{Outils de test}

\begin{table}[H]
  \centering
  \caption{Outils de test utilisés par couche}
  \label{tab:outils_test}
  \begin{tabularx}{\textwidth}{L{4cm}L{3cm}X}
    \toprule
    \thead{Outil} & \thead{Couche} & \thead{Usage} \\
    \midrule
    \rowcolor{lightgray}
    \textbf{PestPHP}       & Backend       & Tests unitaires et d'intégration Laravel \\
    \textbf{Vitest}        & Frontend      & Tests unitaires composants React \\
    \rowcolor{lightgray}
    \textbf{@testing-library/react} & Frontend & Simulation d'interactions utilisateur \\
    \textbf{MSW}           & Frontend      & Mock des appels API dans les tests React \\
    \rowcolor{lightgray}
    \textbf{Postman}       & API REST      & Tests manuels et collections automatisées \\
    \textbf{Newman}        & API REST      & Exécution des collections Postman en CI \\
    \rowcolor{lightgray}
    \textbf{SonarCloud}    & Global        & Analyse statique, couverture, Quality Gate \\
    \textbf{Locust}        & ML / Charge   & Tests de performance du microservice FastAPI \\
    \bottomrule
  \end{tabularx}
\end{table}

% ============================================================

\begin{chapconclusion}
L'étude préliminaire a permis de formaliser l'ensemble des exigences du système :
\textbf{34 user stories} priorisées en MoSCoW, réparties sur \textbf{6 sprints}
pour un total de \textbf{228 story points}. Le périmètre couvre six domaines
fonctionnels : authentification, gestion de projets et équipes, tâches et Kanban,
planification visuelle, intelligence artificielle, et sécurité enterprise avec
journal d'audit analytique. L'architecture distribuée en trois couches
(React~SPA, Laravel~API, FastAPI~ML) garantit l'indépendance de déploiement
et la scalabilité de chaque composant. Les chapitres suivants détaillent
la réalisation sprint par sprint.
\end{chapconclusion}
